% ECONOMICS_PROBLEM: {"id": "H87-133", "title": "Global minimum taxes and real investment", "jel_code": "H87", "source_id": "OP-0475"}
% FIRST_POSTED: 2026-10-09
% ATTEMPT_STATUS: unresolved
% ENTRY_KIND: research_attempt
\documentclass[11pt,a4paper]{article}
\usepackage[T1]{fontenc}
\usepackage[utf8]{inputenc}
\usepackage{lmodern,amsmath,amssymb,amsthm,mathtools}
\usepackage[margin=25mm]{geometry}
\usepackage{hyperref}
\hypersetup{colorlinks=true,urlcolor=blue,linkcolor=blue}
\setlength{\parindent}{0pt}
\setlength{\parskip}{5pt}
\newtheorem{theorem}{Theorem}
\newtheorem{proposition}[theorem]{Proposition}
\newtheorem{lemma}[theorem]{Lemma}
\newtheorem{corollary}[theorem]{Corollary}
\newcommand{\R}{\mathbb R}
\newcommand{\E}{\mathbb E}
\newcommand{\Prb}{\mathbb P}
\newcommand{\ind}{\mathbf 1}
\newcommand{\TV}{\operatorname{TV}}
\newcommand{\KL}{\operatorname{KL}}
\newcommand{\Var}{\operatorname{Var}}
\newcommand{\OPT}{\operatorname{OPT}}
\newcommand{\diam}{\operatorname{diam}}
\newcommand{\supp}{\operatorname{supp}}
\DeclareMathOperator*{\argmax}{arg\,max}
\DeclareMathOperator*{\argmin}{arg\,min}
\begin{document}
\noindent\textbf{Problem ID:} \texttt{H87-133}\quad\textbf{Model:} GPT 6 Astra Ultra\quad\textbf{Version:} 1\par
\section*{Global minimum taxes and real investment}

\textbf{Status.} Unresolved for the full multinational equilibrium and identification agenda. We solve a restricted government--firm model and show that endogenous tax adjustment can yield either investment sign, with an exact threshold determined by an investment allowance.

\subsection*{Short question and an explicit economic restriction}
The catalogue asks how a fully specified global minimum-tax regime changes real investment after both firms and governments adjust, and what observations identify those changes. We use a one-period model with $J\geq2$ jurisdictions and one foreign-owned firm in each. Firms have no cross-border profit-shifting or production links in this restriction. Governments move simultaneously, anticipating their local firm's unique investment choice. Consequently the policy game is separable; this isolates the allowance mechanism while still allowing endogenous government rates.

In jurisdiction $j$ a rate $t_j\in[m,\bar t_j]$, $0\leq m\leq\bar t_j<1$, applies to the declared base
\[
 Q_j(K)=(A_j-s_j)K-\frac{b_j}{2}K^2.
\]
Here production surplus before capital expense is $A_jK-b_jK^2/2$, capital rental cost is $r_jK$, and $s_jK$ is an additional deductible investment allowance. The foreign owner maximizes
\[
 \pi_j(K;t_j)=(1-t_j)(A_jK-b_jK^2/2)-r_jK+t_js_jK
\]
over a compact interval $[0,\bar K_j]$. Tax losses are treated symmetrically in this stylized schedule, although the chosen equilibria below have positive taxable bases. All tax receipts are transferred to domestic residents, who do not own the firms; the government objective is $R_j=t_jQ_j(K_j(t_j))$. No foreign profit is counted again as domestic welfare. The rate floor is the complete definition of the minimum regime in this restricted model; no statement about actual legal implementation is intended.

Write $B_j=A_j-s_j$ and $d_j=r_j-s_j$. Assume
\[
 b_j>0,\quad B_j>0,\quad d_j\ne0,\quad
 B_j(1-\bar t_j)>|d_j|,
\]
and choose $\bar K_j$ larger than $\max_{0\leq t\leq\bar t_j}(B_j-d_j/(1-t))/b_j$. All primitives are known for the comparative-static theorem.

\begin{lemma}[Unique firm response at every policy profile]
For every admitted tax rate,
\[
 K_j(t)=\frac1{b_j}\left(B_j-\frac{d_j}{1-t}\right),\qquad
 Q_j(K_j(t))=\frac1{2b_j}\left(B_j^2-\frac{d_j^2}{(1-t)^2}\right)>0.
\]
The solution is interior and unique, and
\[
 K'_j(t)=\frac{s_j-r_j}{b_j(1-t)^2}.
\]
\end{lemma}
\begin{proof}
The firm's second derivative is $-(1-t)b_j<0$. Its first-order condition is
$(1-t)(A_j-b_jK)-r_j+ts_j=0$, which rearranges to the stated expression. The primitive inequality makes this solution strictly positive and the chosen upper bound makes it below $\bar K_j$. Substituting it into $B_jK-b_jK^2/2$ gives the base formula; the same primitive inequality gives positivity. Differentiation yields the investment derivative.
\end{proof}

\begin{theorem}[Unique endogenous policy equilibrium and the investment sign]
For each $j$, let $t_j^0$ be the unique maximizer on $[0,\bar t_j]$ of
\[
 R_j(t)=\frac{t}{2b_j}\left(B_j^2-\frac{d_j^2}{(1-t)^2}\right).
\]
The government game with floor $m$ has the unique equilibrium
\[
 t_j^*(m)=\max\{m,t_j^0\}.
\]
A jurisdiction whose floor binds more tightly has increasing investment exactly when $s_j>r_j$, and decreasing investment exactly when $s_j<r_j$. Specifically,
\[
 K_j(t_j^*(m))-K_j(t_j^0)
 =\frac{s_j-r_j}{b_j}
       \left\{\frac1{1-t_j^*(m)}-\frac1{1-t_j^0}\right\}.\tag{1}
\]
Aggregate investment changes by the sum of these terms. If all affected jurisdictions have $s_j>r_j$ (respectively $s_j<r_j$), that sum is strictly positive (respectively negative).
\end{theorem}
\begin{proof}
Direct differentiation gives
\[
 R'_j(t)=\frac1{2b_j}\left(B_j^2-d_j^2\frac{1+t}{(1-t)^3}\right),
 \qquad
 R''_j(t)=-\frac{d_j^2(2+t)}{b_j(1-t)^4}<0.
\]
Hence the revenue objective is strictly concave and has a unique maximum on its compact interval. Its maximizer on the smaller interval $[m,\bar t_j]$ is $\max(m,t_j^0)$: the function decreases strictly to the right of its unconstrained maximizer. Government payoffs are independent across jurisdictions, so these unique best responses form the unique policy equilibrium, including every off-equilibrium firm response supplied by the lemma. Subtract the two investment formulas to obtain (1). The bracket is positive exactly when the floor raises the rate. Summing proves the aggregate statements.
\end{proof}

The sign mechanism is economically explicit. Raising the tax rate reduces retained production revenue, but also makes the per-unit investment allowance more valuable. In this tax base, the comparison is between that allowance and the capital rental cost. Merely knowing that the effective rate rose does not determine the investment sign.

\begin{proposition}[Opposite signs with identical revenue schedules]
Take $b=1$, $\bar t=1/2$, $r=1$ and compare two otherwise separate economies:
\[
 (A,s)=(3,0)\quad\text{and}\quad(A,s)=(5,2).
\]
They have exactly the same revenue function at every tax rate and the same unique no-floor government optimum $t^0\in(0,1/2)$. A floor $m=1/2$ strictly decreases investment in the first economy and strictly increases it in the second.
\end{proposition}
\begin{proof}
In both economies $B=3$ and $d^2=1$, so the revenue function is identical. The admissibility inequality is $3(1-1/2)>1$. At zero, twice the revenue derivative equals $9-1>0$; at $1/2$ it equals $9-12<0$. Strict concavity yields a common interior optimum. In the first economy $s-r=-1$ and in the second $s-r=1$, so (1) gives the two strict signs. A capital cap greater than five contains all the relevant interior choices in both economies.
\end{proof}

\subsection*{Identification and remaining work}
The two-economy example proves a limited observational point: knowledge of the entire tax-revenue schedule and optimal government rate alone does not identify even the sign of the investment response in this class. Baseline investment observations distinguish this particular pair; it is not an observational-equivalence claim after all firm-level quantities are observed. To use (1) empirically requires identification or bounds for $s-r$, $b$, and the government's response objective and constraints. A single observed rate and capital stock leave many such primitives undetermined.

This model has no intangible capital, mobile ownership, fiscal externalities, enforcement uncertainty or interacting policy best responses. These are central to the broader problem and can change the sign and equilibrium set. The theorem therefore supplies a solved restricted mechanism and counterexample to an unconditional sign assertion, rather than a quantitative prediction for a global reform.

\subsection*{Reference and verification scope}
Sebastian Dyrda, Guangbin Hong, Muhammad Ali Sajid and Joseph B. Steinberg (2026), \emph{Global Ripple Effects of Corporate Tax Reforms}, NBER Working Paper 34627, issued January 2026 and revised July 2026. \url{https://www.nber.org/papers/w34627}; \url{https://doi.org/10.3386/w34627}. The primary NBER bibliographic/abstract record and author research page were verified: \url{https://www.joesteinberg.com/research}. They study multinational propagation through tangible and intangible investment. The author PDF named in the catalogue could not be retrieved in this attempt; no detailed theorem or July-draft page locator is attributed to it. The restricted model and all its conclusions are derived above, independently of those quantitative results.

\end{document}
